(1) 过阻尼布朗粒子满足 \(dX_t=\sqrt{2D}\,dW_t\)，在区间 \((0,L)\) 两端首次到达即吸收，初始 \(X_0=x\)。由后向方程推导平均首达时间 \(m_1(x)=\mathbb E_x\tau\)，并求其最大值。
(2) 定义 \(F_s(x)=\mathbb E_x e^{-s\tau}\)（\(s>0\)）。写出边值问题并解出双曲函数闭式；由 \(s\to0\) 展开独立复核第一问。
(3) 由扩散核的正弦本征展开求生存概率 \(S(t|x)=P_x(\tau>t)\)；证明 \(m_1(x)=\int_0^\infty S(t|x)\,dt\)，解释长时间指数尾的特征时间。
(4) 利用后向矩递推求 \(m_2(x)=\mathbb E_x\tau^2\) 及中点方差。引入阻力 \(\zeta\) 和温度 \(T\) 时，以 \(D=k_BT/\zeta\) 检验温度依赖与量纲。
